The table below lists the technologies actually used in the Quorum
implementation. Each technology was verified against the repository
(requirements files, configuration, and source code).

\begin{longtable}{@{}p{3.4cm}p{5.4cm}p{6.6cm}@{}}
\toprule
\textbf{Layer} & \textbf{Technology} & \textbf{Usage} \\
\midrule
\endhead
Backend language & Python 3.13 & Main backend, analysis, orchestration \\
Backend framework & FastAPI & REST API and GitHub webhook server \\
ASGI server & Uvicorn & Runs the FastAPI application \\
Data validation & Pydantic & Request/response and model validation \\
Database & PostgreSQL & Persistent application/review/chat data \\
ORM & SQLAlchemy 2.0 & Database access \\
Migrations & Alembic & Database schema migrations \\
GitHub integration & GitHub App, OAuth, Webhooks, REST API & PR events, auth, installation access \\
HTTP client & httpx & GitHub and LLM communication \\
JWT & PyJWT & GitHub App authentication \\
Static analysis & Semgrep (v1.177) & Security evidence over changed files \\
Code analysis & Python standard \texttt{ast} & Structural analysis of changed Python code \\
LLM interface & \texttt{LLMProvider} abstraction & Security, test, and chat reasoning \\
LLM provider (active) & OpenAI (GPT-5.6 Terra / Luna) & Security Agent, Test Writer, Ask Quorum \\
LLM provider (dormant) & Ollama (Qwen2.5-Coder) & Optional local fallback, not in use \\
Test execution & Docker sandbox & Isolated execution of generated tests \\
Sandbox image & \texttt{quorum-sandbox} & pytest + pytest-cov image \\
Backend tests & pytest, pytest-cov & Automated backend tests \\
Frontend & React, TypeScript & Web dashboard \\
Frontend build & Vite & Fast development/build \\
Styling & Tailwind CSS & UI styling \\
UI components & shadcn-style design tokens & Reusable accessible UI components \\
Frontend tests & Vitest, React Testing Library & Component tests \\
Development webhook tunnel & Cloudflare Quick Tunnel & Expose local FastAPI to GitHub \\
\bottomrule
\end{longtable}

\subsection{LLM Configuration}

The LLM is accessed through a common \texttt{LLMProvider} interface so that
no agent depends on a concrete provider. The active provider is OpenAI, with
per-role models:

\begin{itemize}
  \item Security Agent $\rightarrow$ \texttt{gpt-5.6-terra}
  \item Test Writer $\rightarrow$ \texttt{gpt-5.6-terra}
  \item Ask Quorum $\rightarrow$ \texttt{gpt-5.6-luna}
\end{itemize}

A single \texttt{OPENAI\_API\_KEY} authenticates all roles and is stored only
server-side. The \texttt{OllamaProvider} exists in the codebase as an optional
local fallback but is not currently used.